\begin{theorem}[固定分辨率半单不变量的容量因子化]\label{thm:conclusion-window6-semisimple-invariants-factor-through-capacity}
\leanverified{paper\_conclusion\_window6\_semisimple\_invariants\_factor\_through\_capacity}
固定分辨率 \(m\ge 1\)，记
\[
A_m:=\CC[\mathcal R_m^{\mathrm{bin}}],
\qquad
\mathfrak C_m:=\{\mathcal C_m^{\mathrm{cont}}:\ \mathcal C_m^{\mathrm{cont}}\text{ 来自第 }m\text{ 层 bin-fold}\}.
\]
设 \(F\) 为任意取值于某集合 \(\mathcal V\) 且仅依赖有限维半单复 \(\ast\)-代数同构型的函子性不变量。则存在唯一映射
\[
\widetilde F_m:\mathfrak C_m\longrightarrow \mathcal V
\]
使得
\[
F(A_m)=\widetilde F_m(\mathcal C_m^{\mathrm{cont}}).
\]
并且同一条连续容量曲线还唯一决定
\[
\{n_d(m)\}_{d\ge 1},
\qquad
\{S_q(m)\}_{q>0},
\qquad
G_m^{\mathrm{bin}}\cong \prod_{d\ge 1}\mathfrak S_d^{\,n_d(m)}.
\]
\end{theorem}
\begin{proof}
由推论 \ref{cor:conclusion-fold-capacity-wedderburn-exact-bit-threshold}，固定 \(m\) 时连续容量曲线
\[
\mathcal C_m^{\mathrm{cont}}(T)=\sum_{x\in X_m}\min(d_x,T)
\]
唯一决定纤维直方图
\[
n_d(m)=\#\{x\in X_m:\ d_x=d\},
\]
并从而唯一决定 Wedderburn 型
\[
A_m\cong \bigoplus_{d\ge 1}M_d(\CC)^{\oplus n_d(m)}.
\]
由于 \(F\) 只依赖于有限维半单复 \(\ast\)-代数的同构型，它便自动唯一经由 \(\mathfrak C_m\) 因子化，得到所述 \(\widetilde F_m\)。

另一方面，正阶矩谱满足
\[
S_q(m)=\sum_{x\in X_m}d_x^{\,q}=\sum_{d\ge 1}n_d(m)\,d^q
\qquad(q>0),
\]
故也由同一直方图唯一决定。最后，fold-gauge 群的显式分解
\[
G_m^{\mathrm{bin}}
\cong
\prod_{x\in X_m}\mathfrak S_{d_x}
\cong
\prod_{d\ge 1}\mathfrak S_d^{\,n_d(m)}
\]
说明同一数据还唯一锁定群论侧的块型。证毕。
\end{proof}

\begin{corollary}[审计偶数窗上最小退化理想的稳定刚性]\label{cor:conclusion-window6-minideal-stable-rigidity-audited-even}
对审计偶数窗
\[
m\in\{6,8,10,12\},
\]
记最小退化理想
\[
A_{m,\min}:=\bigoplus_{x\in \mathcal S_{m,*}}M_{d_{\min}(m)}(\CC).
\]
则有显式同构
\[
A_{m,\min}\cong M_{F_{m/2}}(\CC)^{\oplus F_m},
\qquad
A_{m,\min}\otimes\mathcal K\cong \mathcal K^{\oplus F_m}.
\]
因此，对任意此类 \(m,n\)，成立
\[
A_{m,\min}\otimes\mathcal K\cong A_{n,\min}\otimes\mathcal K
\quad\Longleftrightarrow\quad
m=n.
\]
\end{corollary}
\begin{proof}
第一式正是定理 \ref{thm:conclusion-foldbin-minideal-fibonacci-primitive-splitting} 的显式分解。对每个 \(r\ge 1\) 有标准同构
\[
M_r(\CC)\otimes\mathcal K\cong \mathcal K,
\]
故
\[
A_{m,\min}\otimes\mathcal K
\cong
\mathcal K^{\oplus F_m}.
\]
由于 Fibonacci 数列严格递增，
\[
\mathcal K^{\oplus F_m}\cong \mathcal K^{\oplus F_n}
\quad\Longleftrightarrow\quad
F_m=F_n
\quad\Longleftrightarrow\quad
m=n.
\]
证毕。
\end{proof}
